\chapter{Epílogo}
\markboth{Epílogo}{Epílogo}

\inicio{V}{eintidós} años más tarde, en la primavera de 2209 ---si es que puede hablarse de primavera en Ceres---\nobreak, Siete subió por última vez al despacho de la doctora Calderón.

Ella ya no daba clases. Pasaba las mañanas en el mismo despacho de siempre, que el Instituto le había dejado por cortesía, con la misma pizarra y, colgada de un clavo en su esquina superior, la llave de bronce del Gremio, ennegrecida por los años. El método del talonario se enseñaba ahora en los manuales elementales. Los estudiantes lo llamaban \emph{el método Calderón}, cosa que a ella le irritaba profundamente.

Siete había cambiado poco. Le habían sustituido dos articulaciones y actualizado el cerebro tres veces, y el Archivo le había nombrado conservador de la sección terrestre, lo cual era, en palabras de Beto Arce ---ya jubilado---\nobreak, «el primer robot con despacho propio de toda la historia de Ceres».

---Doctora ---dijo Siete, después de los saludos---\nobreak, le traigo algo. Lo he encontrado catalogando los fondos antiguos de la Tierra. Llevaba siglo y medio en una caja que nadie había abierto.

Le tendió una lámina. Inés se puso las gafas.

Era un artículo científico, impreso en una tipografía arcaica. Estaba en inglés antiguo, fechado en 2024, publicado en las actas de un congreso llamado Euro-Par. El título decía: \emph{A Fast Wait-Free Solution to Read-Reclaim Races in Reference Counting}. Lo firmaban cuatro personas: Ivo Gabe de Wolff, Daniel Anderson, Gabriele K.\ Keller y Aleksei Seletskiy.

Inés empezó a leer. Al cabo de un párrafo dejó de leer y empezó a pasar páginas, cada vez más deprisa. Siete la observaba en silencio.

En la página cinco había una placa de sesenta y cuatro casillas, con las dieciséis de la izquierda reservadas para un número que los autores llamaban \emph{peso}. En la página seis, un procedimiento para cambiar el cartel: sumar $C$ en el nuevo, trocar, restar los sobrantes en el viejo. En la página ocho, una recarga con un cotejo que salía en cuanto veía más peso del que esperaba. Dos requisitos, que los autores llamaban R-1 y R-2: $C \geq 2P$, y que el peso solo bajara de uno en uno o saltara por encima de $P$. Y, en la página ocho también, una demostración de que el peso nunca era negativo, con una definición de \emph{offence} y un multiconjunto $O$.

Inés se quitó las gafas.

---Ofensa ---dijo---\nobreak. Hasta la palabra.

---Hasta la palabra, doctora. Y hay más. En los ensayos comparan su método con otros cuatro: uno con candados, el contador partido de Williams, el de Folly, y uno de Anderson con centinelas. Los resultados son casi idénticos a los nuestros. Pierden contra Anderson en los mismos escenarios, por la misma razón.

---¿Anderson? ¿El de los centinelas?

---El mismo, doctora. Es uno de los cuatro autores. Compitió contra sí mismo y se ganó en la mitad de los escenarios.

Inés se echó a reír. Fue una risa larga, de las que no había tenido en años, y Siete, que había preparado durante el viaje varios discursos de consuelo para el caso de que la noticia la entristeciera, los borró todos de su memoria.

---Ciento sesenta años ---dijo ella, cuando pudo hablar---\nobreak. Ciento sesenta años tarde. Lo tenían en la sección terrestre de la Bóveda, a tres kilómetros de mi despacho, mientras yo me pasaba la noche en vela con una tiza.

---¿Le molesta, doctora?

Inés miró la llave de bronce colgada de la pizarra.

---Al contrario, Siete. Me tranquiliza. ---Dio unos golpecitos en la lámina---\nobreak. Cuando cuatro personas en la Tierra, y una vieja en Ceres, llegan por separado al mismo talonario, a la misma zona roja, a la misma salida temprana, y hasta a la misma palabra para lo mismo, quiere decir que no nos lo hemos inventado. Que estaba ahí, esperando a que alguien mirara el cartel en vez de la hoja. Lo descubrimos, como se descubre una montaña. Las montañas no son de quien las ve primero.

---El Compendio dirá que no fue usted la primera.

---Que lo diga. Es verdad. ---Le devolvió la lámina---\nobreak. Ponlo en el catálogo, con sus nombres bien escritos. Y añade una nota en la entrada del método: que en Ceres se le llama Calderón por un accidente de la historia, y que sus autores son estos cuatro.

Siete guardó la lámina con un cuidado que habría hecho llorar a un archivero.

---Doctora, una última pregunta. Llevo veintidós años queriendo hacérsela. ¿Por qué nadie lo vio antes, en Ceres? Las dieciséis casillas estaban apagadas a la vista de todos, durante cuarenta años.

Inés se quedó pensando. Por la ventana se veía el Gran Salón, lleno de carteles cuyas casillas de la izquierda bajaban, de uno en uno, sin descanso.

---Porque todos miraban la hoja ---dijo por fin---\nobreak. El problema parecía estar en la hoja: el lector llegaba tarde a la hoja, la hoja llegaba a cero demasiado pronto. Así que todos intentaban arreglar la hoja: protegerla con llaves, retrasar su llegada a cero, vigilar quién se acercaba a ella. Y la solución no estaba en la hoja. Estaba en el cartel, que era lo único que el lector tocaba a tiempo. ---Sonrió---\nobreak. Casi siempre es así, Siete. Cuando un problema no tiene solución donde todo el mundo la busca, suele ser porque está en el sitio al que nadie mira. Y el sitio al que nadie mira suele estar apagado, a la izquierda, esperando.

\bigskip
\begin{center}{\color{oxido}$\ast$\quad$\ast$\quad$\ast$}\end{center}
\bigskip

\begin{documento}[Compendio Universal de Ciencia Archivística · 4.ª ed. · nota añadida en 2209]
\textsc{Método del talonario} (también llamado, impropiamente, \emph{método Calderón}). Técnica para leer y cambiar carteles del Archivo sin esperas, sin llaves y sin acumular cristales sin uso. Descubierta de forma independiente en Ceres en 2187 por I.~Calderón. Por petición expresa de esta, se hace constar que el método fue publicado por primera vez en la Tierra en 2024 por I.~G.~de Wolff, D.~Anderson, G.~K.~Keller y A.~Seletskiy, y que a ellos corresponde la prioridad.
\end{documento}
